%% kapa: exam=1 round=27 year=2016 period=2 form=A subject=accounting subject_ko=회계학 seq=1 number=41
%% source: Q-Net 감정평가사 기출문제 (cst00309) artlSeq=5169402 file=extracted/1008/5169402/제27회 감정평가사 1차 문제/제27회 감정평가사 2교시문제(A형).pdf
%% conversion: kordoc 4.24.51 -> tools/gen1.py (auto); review=proofread-sonnet; flags=none
\begin{question}{41}
재무보고를 위한 개념체계 상 재무제표 요소의 정의 및 인식에 관한 설명으로 옳지 않은 것은?
\begin{choices}
  \item 이익의 측정과 직접 관련된 요소는 수익과 비용이다.
  \item 합리적인 추정을 할 수 없는 경우 해당 항목은 재무상태표나 포괄손익계산서에 인식될 수 없다.
  \item 자산이 갖는 미래경제적효익이란 직접으로 또는 간접으로 미래 현금 및 현금성 자산의 기업에의 유입에 기여하게 될 잠재력을 말한다.
  \item 미래경제적효익의 유입과 유출에 대한 불확실성 정도의 평가는 재무제표를 작성할 때 이용가능한 증거에 기초하여야 한다.
  \item 증여받은 재화는 관련된 지출이 없으므로 자산으로 인식할 수 없다.
\end{choices}
\end{question}
\answer{5}
